\documentclass[12pt,a4paper]{article}

\usepackage[utf8]{inputenc}
\usepackage[slovak]{babel}
\usepackage{geometry}
\usepackage{hyperref}
\usepackage{graphicx}
\usepackage{amsmath}
\usepackage{enumitem}
\usepackage{tocloft}
\usepackage{fancyhdr}
\usepackage{xcolor}
\usepackage{tikz}
\usetikzlibrary{positioning}

\geometry{margin=2.5cm}
\setlength{\headheight}{14.5pt}

% Nastavenie hlavičky a pätičky
\pagestyle{fancy}
\fancyhf{}
\fancyhead[L]{Filtrujúci DNS resolver}
\fancyhead[R]{xbotlo01}
\fancyfoot[C]{\thepage}

\title{\textbf{Projekt do predmetu ISA - Filtrujúci DNS Resolver}}
\author{Filip Botlo - xbotlo01}
\date{2. November 2025}

\begin{document}

\maketitle
\tableofcontents
\newpage

\section{Úvod do problematiky}

DNS (Domain Name System) je hierarchický distribuovaný systém pre preklad doménových mien na IP adresy a naopak. Je základným stavebným kameňom internetu, ktorý umožňuje používateľom pristupovať k webovým stránkam pomocou ľahko zapamätateľných názvov namiesto číselných IP adries.

Cieľom projektu bolo implementovať vlastný DNS resolver, ktorý spracúva DNS dotazy a aplikuje na ne filtrovaciu logiku.

\section{Návrh aplikácie}

\subsection{Architektúra systému}

Aplikácia je navrhnutá ako modulárny systém s jasne oddelenými zodpovednosťami. Používateľ (klient) posiela DNS dotaz napríklad cez nástroje \texttt{dig} alebo \texttt{nslookup}. Dotaz prichádza do DNS resolvera, ktorý ho analyzuje a rozhodne, či je možné ho odoslať na DNS server (napr. Google DNS 8.8.8.8), alebo ho treba blokovať podľa pravidiel filtra.

\begin{figure}[h]
\centering
\begin{tikzpicture}[node distance=1.5cm, >=stealth, thick]
\tikzstyle{block} = [rectangle, draw, rounded corners,
                    text centered, minimum height=1cm,
                    minimum width=3cm]
\node[block] (client) {DNS Client \\ (dig, nslookup)};
\node[block, right=of client] (resolver) {DNS Resolver};
\node[block, right=of resolver] (upstream) {Upstream DNS \\ (8.8.8.8)};
\node[block, below=of resolver, yshift=-0.5cm] (filter) {Filter Engine \\ (Block / Allow)};
\draw[->] (client) -- (resolver);
\draw[->] (resolver) -- (upstream);
\draw[->] (resolver) -- (filter);
\end{tikzpicture}
\caption{Architektúra systému}
\end{figure}

\section{Implementácia}

\subsection{Prehľad tried a ich funkcionality}

Aplikácia je implementovaná v jazyku C++ s využitím objektovo orientovaného programovania. Hlavné triedy sú navrhnuté tak, aby každá mala jasne definovanú zodpovednosť a minimalizovala závislosti medzi modulmi.

\subsubsection{Trieda CLIParser}

Trieda je zodpovedná za spracovanie argumentov príkazového riadka. Obsahuje metódy pre parsovanie jednotlivých parametrov, validáciu vstupných hodnôt a generovanie nápovedy. Hlavné funkcie tejto triedy zahŕňajú:

\begin{itemize}
    \item \texttt{parse()} -- hlavná metóda pre spracovanie argumentov
    \item \texttt{validateArguments()} -- kontrola platnosti zadaných parametrov
    \item \texttt{printHelp()} -- zobrazenie nápovedy pre používateľa
\end{itemize}

Trieda podporuje všetky potrebné parametre ako port, upstream DNS server, cestu k filtrovaciemu súboru, verbose mód a štatistiky. Implementuje error handling pre neplatné vstupy.

\subsubsection{Trieda DNSResolver}

Trieda predstavuje jadro aplikácie a koordinuje všetky operácie. Je zodpovedná za vytvorenie sieťových socketov, implementáciu serverového cyklu a správu komunikácie medzi klientmi a upstream DNS serverom. Kľúčové metódy tejto triedy:

\begin{itemize}
    \item \texttt{createClientSocket()} -- vytvorenie socketu pre prijímanie dotazov od klientov
    \item \texttt{createResolverSocket()} -- vytvorenie socketu pre komunikáciu s upstream DNS
    \item \texttt{loadFilterFile()} -- načítanie zoznamu blokovaných domén zo súboru
    \item \texttt{handleQuery()} -- spracovanie jednotlivého DNS dotazu
    \item \texttt{run()} -- hlavný serverový cyklus
\end{itemize}

Serverový cyklus používa \texttt{select()} s \texttt{recvfrom()} pre prijímanie dotazov a \texttt{sendto()} pre odosielanie odpovedí.

\subsubsection{Trieda DNSProtocol}

Trieda \texttt{DNSProtocol} obsahuje všetku logiku súvisiacu s DNS protokolom podľa RFC 1035. Je implementovaná ako statická trieda s utility metódami pre prácu s DNS správami. Hlavné funkcie:

\begin{itemize}
    \item \texttt{parseDomain()} -- parsovanie doménových mien z DNS správ
    \item \texttt{isQuery()} -- kontrola, či je správa DNS dotazom
    \item \texttt{processQuery()} -- hlavná logika filtrovania a spracovania dotazov
    \item \texttt{sendErrorResponse()} -- generovanie chybových DNS odpovedí
\end{itemize}

\subsection{DNS protokol implementácia}

Parsovanie doménových mien je jednou z najzložitejších častí implementácie. Domény sú uložené ako sekvencia labelov, kde každý label má svoju dĺžku nasledovanú obsahom. Resolver dokáže správne dekódovať tieto štruktúry a rekonštruovať plné doménové meno.

\subsection{Filtrovacia logika}

Najskôr sa doména normalizuje -- prevedie sa na malé písmená a odstránia sa nadbytočné znaky ako koncové bodky. Následne sa kontroluje presná zhoda s blokovanými doménami. Pre podporu subdomén sa implementuje algoritmus, ktorý kontroluje, či normalizovaná doména končí bodkou nasledovanou niektorou z blokovaných domén. Toto umožňuje blokovať všetky poddomény bez potreby explicitne uvádzať každú z nich. Filtrovacia logika podporuje len A záznamy (typ 1), pričom ostatné typy záznamov sú odmietnuté s chybovým kódom NOTIMP.

\subsection{Rozšírenia aplikácie}

\subsubsection{Verbose mód}

Verbose mód je implementovaný ako voliteľný parameter \texttt{-v} alebo \texttt{--verbose}. Pri aktivácii tohto módu resolver vypisuje podrobné informácie o každom spracovanom DNS dotaze, čo značne uľahčuje ladenie a monitorovanie činnosti aplikácie. Implementácia verbose módu je integrovaná do triedy \texttt{DNSResolver} a \texttt{DNSProtocol}, kde sa kontroluje flag \texttt{verbose\_} pred každým výpisom diagnostických informácií.

\subsubsection{Štatistiky}

Rozšírenie štatistík je implementované ako voliteľný parameter \texttt{--stats}. Toto rozšírenie umožňuje zbieranie a zobrazenie štatistík o činnosti DNS resolvera na konci jeho behu. Štatistiky zahŕňajú nasledujúce metriky:

\begin{itemize}
    \item Celkový počet dotazov
    \item Blokované dotazy
    \item Preposlané dotazy
    \item Ostatné typy dotazov (NOTIMP)
    \item Presné zhody
    \item Zhody s poddomenami
\end{itemize}

Implementácia štatistík využíva štruktúru \texttt{Statistics} definovanú v \texttt{DNSProtocol.hpp}.

\section{Testovanie}

Testovanie DNS resolvera bolo navrhnuté tak, aby pokrylo všetky kritické aspekty implementácie a overilo správne správanie systému podľa očakávaní. Testovanie sa skladalo z manuálnych testov pomocou nástroja \texttt{dig} a automatizovaných testov pomocou Python skriptu.

\subsection{Manuálne testovanie pomocou dig}

Nástroj \texttt{dig} bol zvolený ako hlavný nástroj pre testovanie, pretože ide o štandardnú implementáciu DNS klienta, ktorá správne generuje DNS dotazy podľa RFC 1035. Pre každý test som definoval jasné očakávania na základe implementácie resolvera.

\subsubsection{Testovanie A záznamov -- povolené domény}

Pri testovaní povolených domén som očakával, že resolver prijme dotaz a správne ho preposiela na upstream DNS server, zachová pôvodné Query ID, vráti správnu odpoveď so statusom \texttt{NOERROR} a obsah odpovede bude obsahovať IP adresu z upstream DNS servera.

\textbf{Príklad testu:}

\begin{verbatim}
dig @127.0.0.1 -p 5353 google.com
\end{verbatim}

Pomocou tohto dotazu som overil, že doména \texttt{google.com}, ktorá nie je v zozname blokovaných domén, je správne preposlaná na resolver \texttt{8.8.8.8}. Očakával som návrat správnej IP adresy a status kódu \texttt{NOERROR}.

\subsubsection{Testovanie A záznamov -- blokované domény}

Pri testovaní blokovaných domén som očakával, že resolver rozpozná doménu v zozname blokovaných domén (presná zhoda alebo subdoména), nezasiela dotaz na upstream DNS server, vytvorí chybovú odpoveď so statusom \texttt{REFUSED} (kód 5) a zachová Query ID z pôvodného dotazu.

\textbf{Príklad testu:}

\begin{verbatim}
dig @127.0.0.1 -p 5353 facebook.com
\end{verbatim}

Doména \texttt{facebook.com} je explicitne uvedená v súbore \texttt{blocked\_domains.txt}, preto som očakával, že dotaz bude odmietnutý s kódom \texttt{REFUSED}.

\subsubsection{Testovanie subdomén}

Kľúčová funkčnosť resolvera je schopnosť blokovať aj subdomény blokovaných domén. Pre test som použil subdoménu \texttt{www.facebook.com}, pričom som očakával, že bude blokovaná, pretože \texttt{facebook.com} je v zozname blokovaných domén.

\begin{verbatim}
dig @127.0.0.1 -p 5353 www.facebook.com
\end{verbatim}

Očakával som status kód \texttt{REFUSED}, čo potvrdilo, že implementácia správne kontroluje, či dotazovaná doména končí bodkou nasledovanou niektorou z blokovaných domén.

\subsubsection{Testovanie rôznych typov DNS záznamov}

Podľa implementácie resolver podporuje len A záznamy (typ 1). Pre ostatné typy záznamov som očakával status kód \texttt{NOTIMP} (kód 4).

\textbf{Test MX záznamov:}

\begin{verbatim}
dig @127.0.0.1 -p 5353 google.com MX
\end{verbatim}

Očakával som status \texttt{NOTIMP}, pretože MX záznamy (typ 15) nie sú podporované. Rovnaké očakávanie som mal aj pre AAAA záznamy (typ 28), NS, TXT, CNAME a ďalšie typy záznamov.

\subsubsection{Testovanie zachovania Query ID}

Dôležitým aspektom DNS protokolu je zachovanie Query ID v odpovedi. V aplikácii nemanipulujem s Query ID, preto som očakával, že pôvodné Query ID z dotazu bude zachované v odpovedi, čo je kritické pre správne spárovanie odpovede s pôvodným dotazom klientom.

\subsection{Automatizované testovanie}

Pre komplexnejšie testovanie som vytvoril Python testovací skript \texttt{simple\_test.py}, ktorý automatizuje opakujúce sa testy a umožňuje systematické overenie funkcionality.

\subsubsection{Python testovací skript}

Testovací skript \texttt{simple\_test.py} pokrýva nasledujúce kategórie testov:

\begin{itemize}
    \item \textbf{Kompilácia a základné funkcie} -- overenie úspešnej kompilácie a zobrazenia nápovedy
    \item \textbf{Validácia argumentov} -- testovanie rôznych kombinácií argumentov, poradia parametrov, neplatných hodnôt
    \item \textbf{IPv4/IPv6 podpora} -- testovanie s rôznymi typmi adries a doménovými menami DNS serverov
    \item \textbf{Spracovanie súborov} -- testovanie rôznych formátov koncov riadkov (Linux, Windows, Mac), komentárov, prázdnych riadkov
    \item \textbf{DNS dotazy} -- systematické testovanie blokovaných a povolených domén, subdomén, case sensitivity
    \item \textbf{Rôzne typy dotazov} -- overenie správneho odmietania nepodporovaných typov (MX, AAAA, TXT, NS, CNAME, atď.)
    \item \textbf{Chybové scenáre} -- testovanie neexistujúcich súborov, neplatných DNS serverov, timeoutov
\end{itemize}

Skript automaticky spúšťa DNS resolver v pozadí, vykonáva testy a následne vyhodnotí výsledky. Poskytuje detailné štatistiky o úspešnosti testov a identifikuje zlyhané testy s popisom chýb.

\subsection{Výsledky testov}

Všetky testy boli navrhnuté na základe konkrétnych očakávaní založených na implementácii resolvera. Testovanie potvrdilo:

\begin{itemize}
    \item Správne filtrovanie blokovaných domén (presné zhody aj subdomény)
    \item Správne preposielanie povolených domén na upstream DNS server
    \item Zachovanie Query ID v odpovediach
    \item Správne odmietanie nepodporovaných typov záznamov (\texttt{NOTIMP})
    \item Podpora IPv4 aj IPv6 klientov
    \item Správne spracovanie rôznych formátov filtrovacieho súboru
    \item Case-insensitive porovnávanie domén
\end{itemize}

Automatizovaný testovací skript úspešne prešiel všetky testy, čo potvrdilo správnu implementáciu všetkých požadovaných funkcií.

\section{Záver}

Vďaka projektu som sa naučil pracovať s DNS protokolom na nízkej úrovni a lepšie pochopil jeho štruktúru podľa RFC 1035. Získal som praktické skúsenosti so sieťovým programovaním v C++ a s implementáciou vlastného resolvera. Projekt úspešne implementuje filtrujúci DNS resolver s podporou IPv4/IPv6, filtrovania domén a štatistík. Všetky testy prešli úspešne.

\section{Literatúra a zdroje}

\begin{enumerate}
    \item \url{https://tools.ietf.org/html/rfc1035}
    \item \url{https://man7.org/linux/man-pages/man2/socket.2.html}
    \item \url{https://man7.org/linux/man-pages/man2/sendto.2.html}
    \item \url{https://man7.org/linux/man-pages/man2/recvfrom.2.html}
    \item \url{https://man7.org/linux/man-pages/man2/select.2.html}
    \item \url{https://man7.org/linux/man-pages/man7/ipv6.7.html}
\end{enumerate}

\end{document}

